% Edit the title, individual panel descriptions, and shared notes here.
% All figure placements read these commands; edit once to update every continuation.
\newcommand{\FigTwoTitle}{Pseudotime organization and learned RF transition graph characteristics.}
\newcommand{\FigTwoA}{Normalized Monocle 3 pseudotime in Walesky hepatocytes sampled at 0~h (n = 3,401), 6~h (419), 24~h (1,043), 48~h (910), and 96~h (809). White circles indicate medians and thick black segments span the interquartile range; pseudotime is dimensionless.}
\newcommand{\FigTwoB}{Area under the receiver operating characteristic curve (AUROC) from five-fold stratified cross-validation of the four scale-specific RF transition classifiers. Points are individual folds, bars show mean AUROC, and diamonds show out-of-bag (OOB) accuracy. Black error bars show mean AUROC plus or minus 1.96 times the fold SD divided by the square root of five; they summarize variation among overlapping cross-validation folds, rather than independent biological replicates. The classifier training sets contain 153,045, 18,855, 46,935, and 40,950 labeled cell pairs for Scales 1--4, respectively.}
\newcommand{\FigTwoC}{Mean row-wise Pearson correlation between learned RF transition probabilities and distance-only Gaussian transition probabilities, calculated over source cells in each temporal scale.}
\newcommand{\FigTwoD}{Mean Shannon entropy of the outgoing RF transition probability distribution, in nats, across 3,401, 419, 1,043, and 910 source cells in Scales 1--4, respectively. These are descriptive graph summaries.}
\newcommand{\FigTwoNotes}{Scales 1--4 connect 0--6, 6--24, 24--48, and 48--96~h, respectively. RF, Random Forest. Raw probabilities below 0.30 are removed before row normalization, with the 15 highest-probability targets retained when no target meets that threshold; normalized probabilities below 0.001 are subsequently removed.}
